\chapter{Đọc thêm 3: Các vấn đề chất lượng dữ liệu đối với nhà nghiên cứu kinh doanh và tài chính}
\label{ch:M2L4-reading3}

% Nguồn: Liu, Grace. "Data Quality Problems Troubling Business and Financial Researchers: A Literature Review and Synthetic Analysis." \emph{Journal of Business \& Finance Librarianship}, 2020. https://digitalcommons.wcupa.edu/lib_facpub/13/
% Phần cần đọc: tr. 29--36 (~25 phút)
% Nội dung bắt buộc đọc: được kiểm tra trong mọi bài đánh giá (kể cả Quiz).

\section{Thông tin nguồn}

\begin{itemize}
  \item \textbf{Nguồn:} Liu, Grace. "Data Quality Problems Troubling Business and Financial Researchers: A Literature Review and Synthetic Analysis." \emph{Journal of Business \& Finance Librarianship}, 2020. \url{https://digitalcommons.wcupa.edu/lib_facpub/13/}
  \item \textbf{Phần cần đọc:} tr. 29--36 (~25 phút)
\end{itemize}

câu hỏi nghiên cứu, cung cấp cho các nhà nghiên cứu kinh doanh và thủ thư thông tin về các vấn đề dữ liệu phổ biến và thảo luận về hàm ý của chúng đối với tham khảo kinh doanh và tư vấn nghiên cứu.

\section{Các Vấn đề Chất lượng Dữ liệu (Data Quality Problems) Phổ biến}

\section{Giá trị bị thiếu (Missing Values)}

Giá trị bị thiếu là một trong những vấn đề chất lượng dữ liệu phổ biến nhất. Các tệp dữ liệu CRSP được phát hiện bị thiếu lợi suất hủy niêm yết (delisting returns), lợi suất quỹ tương hỗ (mutual fund returns) và các danh mục quỹ tương hỗ bắt buộc (Shumway, 1997; Elton et al., 2001; Wisen, 2002; Schwarz \& Potter, 2016). Một lượng lớn dữ liệu bị bỏ sót đã được xác định trong Compustat (Boritz \& No, 2013; Casey et al., 2016; Heitzman \& Lester, 2020). Sự bỏ sót dữ liệu về ``chủ sở hữu cuối cùng toàn cầu'' (global ultimate owner) đã được báo cáo trong Orbis (Jaraitė et al., 2013). Dữ liệu hạn mức tín dụng (credit line) và dữ liệu nợ chưa thanh toán (outstanding debt) trong Capital IQ được quan sát thấy có một lượng lớn các giá trị bị thiếu (Mathers \& Giacomini, 2016; Lee, 2017). Dữ liệu bảng cân đối kế toán, dữ liệu số lượng nhân viên được báo cáo và dữ liệu lợi suất thực tế trên tài sản chương trình hưu trí (actual return on plan assets) được phát hiện bị thiếu trong Worldscope (Weiß \& Mühlnickel, 2014; McGuire et al., 2016; Nobes \& Stadler, 2018). Thông tin cho những người dự thầu, điều khoản chấm dứt (termination provisions), thông báo mua lại cổ phần (share repurchase announcements), thông báo tiếp quản (takeover announcements) và các sự kiện liên quan đến sáp nhập được phát hiện bị thiếu trong SDC (Faccio \& Masulis, 2005; Boone \& Mulherin, 2007; Banyi et al., 2008; Mulherin \& Aziz Simsir, 2015). Dữ liệu phát hành trái phiếu, dữ liệu tài chính cổ phiếu, dữ liệu đối thủ cạnh tranh và dữ liệu cây sản phẩm trong Mergent Online được phát hiện không đầy đủ (Tallapally, 2009; Ma et al., 2011; Berrios, 2013; Lu \& Shang, 2017).

Các nhà nghiên cứu thường xác định các giá trị bị thiếu bằng cách so sánh trực tiếp dữ liệu với dữ liệu gốc trong các hồ sơ gửi lên SEC và gần đây hơn là với dữ liệu tương tác được gắn thẻ XBRL (XBRL-tagged) từ hệ thống EDGAR của SEC (Tallapally et al., 2011; Boritz \& No, 2013; Chychyla \& Kogan, 2015; Schwarz \& Potter, 2016). Giá trị bị thiếu xảy ra phổ biến hơn đối với các khái niệm kế toán phức tạp như ``lỗ hoạt động ròng'' (net operating loss), ``hạn mức tín dụng'', ``nợ chưa thanh toán'', ``lợi suất thực tế trên tài sản chương trình hưu trí'' (Mathers \& Giacomini, 2016; Lee, 2017; Nobes \& Stadler, 2018; Heitzman \& Lester, 2020). Giá trị bị thiếu cũng thường xuyên xảy ra hơn đối với các giao dịch phức tạp cần nhiều nỗ lực để theo dõi các thay đổi như mua lại cổ phần và sáp nhập và mua lại (M\&A) (Banyi et al., 2008; Chapman \& Klein, 2009; Netter et al., 2011). Giá trị bị thiếu thường được quan sát thấy nhiều hơn đối với dữ liệu trong các phần chú thích hoặc dữ liệu không xuất hiện trực tiếp trên báo cáo tài chính, chẳng hạn như ``lợi suất thực tế trên tài sản chương trình hưu trí'' hoặc ``tín dụng tuần hoàn chưa rút'' (undrawn revolving credit) (Mathers \& Giacomini, 2016; Nobes \& Stadler, 2018). Trong Orbis, giá trị bị thiếu cũng có thể xảy ra do giới hạn về lượng dữ liệu được phép tải xuống (Kalemli-Ozcan et al., 2019).

Các nhà nghiên cứu đôi khi áp dụng các quy trình hoặc bộ lọc đặc biệt để loại trừ các giá trị bị thiếu khỏi mẫu nghiên cứu. Tuy nhiên, thực tiễn này chắc chắn có thể tạo ra thiên lệch bỏ sót (omission bias) hoặc thiên lệch chọn mẫu (selection biases) (Elton et al., 2001; Weiß \& Mühlnickel, 2014). Việc loại bỏ tất cả các quan sát chứa giá trị bị thiếu là một chiến lược ngây thơ và có thể có tác động rõ rệt đến sức mạnh thống kê của các phép kiểm định (Hribar, 2016, tr. 63). Loại trừ các giá trị bị thiếu có thể tạo ra các kết quả sai lệch và số lượng lớn các giá trị bị thiếu có thể làm cho một cơ sở dữ liệu không thể sử dụng được cho nghiên cứu cụ thể (Francis et al., 2016; Lee, 2017). Việc thiếu lợi suất hủy niêm yết có thể dẫn đến thiên lệch hủy niêm yết (delisting bias) và các thiên lệch dữ liệu chưa biết khác làm nhiễu kết quả thực nghiệm (Shumway, 1997; Shumway \& Warther, 1999). Các giá trị bị thiếu có thể được giảm bớt bằng cách so sánh các tập dữ liệu với hồ sơ SEC hoặc các nguồn dữ liệu khác. Casey et al. (2016) đã đề xuất một

% [PDF p.30 | in p.29]
Mô hình Cân bằng Báo cáo Tài chính được Điều chỉnh (Modified Financial Statement Balancing Model) để giải quyết một phần các vấn đề về giá trị bị thiếu hoặc các mục nhập sai sót và khôi phục chúng thành các điểm dữ liệu có thể sử dụng được như số 0 hoặc các khoản tổng hợp.

\section{Lỗi Dữ liệu (Data Errors)}

Lỗi dữ liệu là một vấn đề chất lượng dữ liệu phổ biến khác. Nhiều lỗi số học, lỗi mã hóa, lỗi ngày sáp nhập, lỗi vị thế danh mục đầu tư đã được tìm thấy trong CRSP (Courtenay \& Keller, 1994; Elton et al., 2001; Schwarz \& Potter, 2016). Các lỗi đánh máy hoặc lỗi làm tròn, lỗi phân loại, lỗi tính toán, các biến số kiểm toán viên bị mã hóa sai đã được tìm thấy trong Compustat (San Miguel, 1977; Chychyla \& Kogan, 2015); sự phân loại sai các thực thể và cổ đông đã được tìm thấy trong Orbis (Monasterolo et al., 2017). Sự chênh lệch (disparity) giữa dữ liệu hạn mức tín dụng Capital IQ và dữ liệu trong hồ sơ 10-K được phát hiện là đáng kể (Mathers \& Giacomini, 2016). Mã nhà phân tích và ngày công bố thu nhập trong I/B/E/S được quan sát thấy bị lỗi báo cáo (Acker \& Duck, 2009; Roger, 2017). Các lỗi dữ liệu về khối lượng (volume), giá cả, cổ phiếu, chỉ số lợi suất và tổng lợi suất, phân loại, ngày tháng và thông tin hủy niêm yết đã được tìm thấy trong Datastream (Bloom et al., 2004; Rossi, 2011; Brückner, 2013). Các lỗi trong đơn vị chào bán (unit offers), phân loại ngành, doanh số, tài sản, giá chào bán, giá thị trường hoặc giá trị thương vụ đã được báo cáo trong SDC. Tỷ lệ lỗi cao đối với số cổ phiếu đang lưu hành sau khi phát hành (post-issue shares outstanding), số lượng cổ phiếu phân bổ vượt mức (overallotment shares) được thực thi và số lượng nhà bảo lãnh phát hành quản lý (managing underwriters) cũng được báo cáo trong SDC (Ritter, 2019).

Nhiều nhà nghiên cứu xác định lỗi dữ liệu bằng cách so sánh các cơ sở dữ liệu khác nhau hoặc bằng cách so sánh cơ sở dữ liệu với dữ liệu thu thập thủ công của họ (Beedles \& Simkowitz, 1978; Courtenay \& Keller, 1994). Lỗi dữ liệu có thể xảy ra ở các dạng đơn giản như lỗi đánh máy hoặc lỗi làm tròn. Nhưng thường xuyên hơn, lỗi dữ liệu xảy ra ở các dạng phức tạp hơn. Lỗi dữ liệu có thể là kết quả của việc bao gồm hoặc loại trừ không đúng một số khoản mục kế toán trong các phép tính toán, chẳng hạn như gộp cả nghiên cứu theo hợp đồng (contract research) vào chi phí R\&D, xử lý sai các hoạt động kinh doanh và đầu tư, loại trừ các khoản chênh lệch dồn tích phải trả (accrued imbalances payable) [?] khỏi các khoản phải trả (San Miguel, 1977; Shi \& Zhang, 2011). Các lỗi có nhiều khả năng xảy ra đối với các khái niệm tài chính phức tạp như ``giá vốn hàng bán'' (cost of goods sold), ``lợi nhuận gộp'' (gross profit) hoặc ``lỗ hoạt động ròng'' hơn là đối với các khái niệm đơn giản như ``tổng tài sản'', ``tổng nợ phải trả'' hoặc ``thu nhập ròng'' (Chychyla \& Kogan, 2015; Heitzman \& Lester, 2020). Lỗi thường xuyên xảy ra hơn đối với các giao dịch phức tạp như sáp nhập và mua lại, thay đổi sàn giao dịch, hủy niêm yết hoặc chia tách cổ phiếu (Andrikopoulos et al., 2007; Rossi, 2011; Nobes \& Stadler, 2018; Ritter, 2019). Lỗi mã hóa sai trong các biến số kiểm toán viên xảy ra thường xuyên hơn khi có sự thay đổi kiểm toán viên (Utke, 2018). Tỷ lệ lỗi nói chung cao hơn đối với các khoản mục được báo cáo trong phần chú thích so với các khoản mục được báo cáo trên báo cáo kết quả hoạt động kinh doanh hoặc bảng cân đối kế toán (Kinney \& Swanson, 1993; Bratten et al., 2016). Lỗi phân loại (classification errors) có thể xảy ra do sự hiểu nhầm các khái niệm kinh doanh và tài chính như cổ phiếu phổ thông hoặc phân loại ngành (Ince \& Porter, 2006; Ritter, 2019). Về các công ty nước ngoài hoặc giao dịch nước ngoài, lỗi có thể phát sinh từ việc dịch hoặc diễn giải không đúng các thuật ngữ kế toán nước ngoài hoặc các điều kiện của công ty nước ngoài (Brückner, 2013; Nobes \& Stadler, 2018). Các lỗi cũng phát sinh khi cơ sở dữ liệu không cập nhật dữ liệu kịp thời khi có những thay đổi về sàn giao dịch, phân loại ngành, kiểm toán viên, sự trình bày lại (restatement), v.v. (Rossi, 2011; Chychyla \& Kogan, 2015; Utke, 2018).

Lỗi dữ liệu có thể làm bóp méo kết quả của các nghiên cứu liên quan theo nhiều cách (Acker \& Duck, 2009). Việc làm tròn giá đến đồng xu (penny) gần nhất có thể không phải là một lỗi nghiêm trọng (hard error), nhưng nó có thể gây ra những khác biệt không hề nhỏ trong lợi suất tính toán khi giá ở mức nhỏ (Ince \& Porter, 2006). Các lỗi báo cáo trong mã nhà phân tích có thể tác động đến việc đánh giá

% [PDF p.31 | in p.30]
đặc điểm của các nhà phân tích và có thể làm thiên lệch các nghiên cứu thực nghiệm dựa vào việc theo dõi các nhà phân tích (Roger, 2017). Những lỗi lớn có thể ảnh hưởng đến một số thuộc tính của mẫu ở một mức độ không tương xứng với số lượng nhỏ của chúng, đưa vào các thiên lệch và làm nhiễu loạn các phân tích thống kê (Rosenberg \& Houglet, 1974). Một lượng lớn thông tin không chính xác sẽ khiến các nhà nghiên cứu cần xác minh dữ liệu một cách độc lập, điều này làm suy yếu giá trị của việc sử dụng các cơ sở dữ liệu thương mại (Elton et al., 2001). Lỗi dữ liệu có thể được phát hiện và giảm bớt bằng cách kiểm tra chéo (cross-checking) với các nguồn khác. Mô hình Cân bằng Báo cáo Tài chính được Điều chỉnh có thể được sử dụng để tìm các mục nhập sai sót (Casey et al., 2016). Các phương pháp lấy mẫu thống kê thường được sử dụng để xác định lỗi dữ liệu và các giá trị ngoại lai (CRSP LLC, 2020c).

\section{Sự chênh lệch (Discrepancies)}

Nhiều nhà nghiên cứu đã tìm thấy các sự chênh lệch khi so sánh các cơ sở dữ liệu hoặc tập dữ liệu khác nhau (Grinblatt et al., 1984, Sarig \& Warga, 1989; Guenther \& Rosman, 1994; Courtenay \& Keller, 1994; Kern \& Morris, 1994; Kahle \& Walkling, 1996; Elton et al., 2001; Yang et al., 2003; Ulbricht \& Weiner, 2005; Daske et al., 2013; Bollaert \& Delanghe, 2015; McGuire et al., 2016; Tobek \& Hronec, 2018). Sự chênh lệch có thể là do sự khác biệt về độ bao phủ cơ sở dữ liệu, định nghĩa, chính sách mã hóa, mã định danh, phân loại, mô hình tính toán, thiên lệch chọn mẫu, hoặc lỗi dữ liệu (Sarig \& Warga, 1989; Courtenay \& Keller, 1994; Kern \& Morris, 1994; Yang et al., 2003; Tallapally, 2009; Dreyer \& Hines, 2014; Nam, et al., 2017).

Thường xuyên nhất, sự chênh lệch được phát hiện bằng cách so sánh hai cơ sở dữ liệu hoặc tập dữ liệu khác nhau. Tuy nhiên, sự chênh lệch cũng có thể xảy ra trong phạm vi một cơ sở dữ liệu, đặc biệt là trong các cơ sở dữ liệu tập hợp (aggregator databases) mà thu thập dữ liệu từ các nguồn khác nhau. Ví dụ, các nhà nghiên cứu đã phát hiện ra rằng EPS thực tế trong I/B/E/S khác với EPS thực tế được nhà phân tích suy ra ở 39\% số lần (Brown \& Larocque, 2011, 2013). Sự chênh lệch phổ biến hơn đối với các giao dịch phức tạp như chia tách cổ phiếu, cổ tức bằng cổ phiếu, sáp nhập và mua lại, những thay đổi kế toán hoặc các hoạt động bị ngừng lại (Grinblatt et al., 1984; Courtenay \& Keller, 1994, Kern \& Morris, 1994). Sự chênh lệch là phổ biến hơn đối với các khoản mục dữ liệu dựa trên sự gán chủ quan như mã SIC (Guenther \& Rosman, 1994; Kahle \& Walkling, 1996; Stasch, 2014); Sự chênh lệch có thể xảy ra do các quy trình chuẩn hóa trong cơ sở dữ liệu có thể không nhất thiết mang lại lợi ích cho nghiên cứu thực nghiệm (Tallapally et al., 2012; Chychyla \& Kogan, 2014). Mức độ chênh lệch cao hơn giữa các cơ sở dữ liệu được xác định đối với dữ liệu quốc tế, đặc biệt là dữ liệu về các quốc gia đang phát triển (Lara et al., 2006; Daske et al., 2013; McGuire et al., 2016; Hines \& Sharma, 2018). Sự chênh lệch lớn về tính khả dụng của dữ liệu giữa các quốc gia khác nhau phần lớn là do các thông lệ kế toán và quy tắc nộp hồ sơ khác nhau dựa trên các tập quán lịch sử và văn hóa khác nhau (Green, 2007). Firat (2002) đã xác định ít nhất ba loại tính không đồng nhất (heterogeneities) (tính không đồng nhất ở cấp độ dữ liệu, về mặt bản thể học và thời gian) giữa các cơ sở dữ liệu. Sự khác biệt phát sinh khi các cơ sở dữ liệu chọn các đơn vị, tỷ lệ hoặc định dạng khác nhau để trình bày dữ liệu của cùng một công ty; hoặc khi chúng áp dụng các định nghĩa kế toán, chính sách chuyển đổi tiền tệ khác nhau; hoặc khi các giá trị hay định nghĩa thực thể thuộc về các thời điểm, hoặc khoảng thời gian khác nhau (Firat, 2002).

Sự chênh lệch trên các cơ sở dữ liệu hàm ý rằng sẽ có rủi ro khi dựa vào bất kỳ cơ sở dữ liệu đơn lẻ nào (McGuire, 2016). Sự chênh lệch có thể dẫn đến ``hiệu ứng cơ sở dữ liệu'' (database effect), có nghĩa là các nhà nghiên cứu sẽ đi đến các kết luận khác nhau dựa trên các cơ sở dữ liệu khác nhau. Những khác biệt trong cơ sở dữ liệu có thể ảnh hưởng đến các lựa chọn mẫu, các suy diễn

% [PDF p.32 | in p.31]
về tổng thể, và có thể ảnh hưởng sâu hơn tới kết quả của nghiên cứu thực nghiệm (Kern \& Morris, 1994; Kahle \& Walkling, 1996; Lara et al., 2006; McGuire et al., 2016). Nhiều nhà nghiên cứu chọn sử dụng nhiều nguồn dữ liệu cho nghiên cứu của họ; tuy nhiên, cần phải nỗ lực lớn để so sánh các chính sách mã hóa, đối chiếu (matching), lọc, và làm sạch dữ liệu (Chakrabarty \& Trzcinka, 2006).

\section{Các Thiên lệch (Biases)}

Những lo ngại về các thiên lệch trong cơ sở dữ liệu được thảo luận rộng rãi. Thiên lệch phóng đại (upward bias), thiên lệch hủy niêm yết, thiên lệch bỏ sót, thiên lệch sống sót (survivorship bias), thiên lệch ươm tạo (incubation bias), thiên lệch điền lùi (backfill bias) (hoặc thiên lệch 'lịch sử tức thời' - 'instant-history' bias), và thiên lệch trùng lặp (duplication bias) đã được phát hiện trong CRSP (Loughran \& Ritter, 1995; Shumway, 1997; Canina et al., 1998; Shumway \& Warther, 1999; Elton et al., 2001; Wisen, 2002; Yan, 2007; Evans, 2010; Jorion and Schwarz, 2017; CRSP, 2020). Thiên lệch sống sót, thiên lệch chọn lọc sau sự kiện (ex-post-selection bias), và thiên lệch nhìn trước (look-ahead bias) đã được tìm thấy trong Compustat (Ball, 1979; Banz \& Breen, 1986). Thiên lệch sống sót và thiên lệch chọn mẫu đã được báo cáo trong Datastream và Orbis (Ince \& Porter, 2006; Andrikopoulos et al., 2007; Kalemli-Ozcan et al., 2019).

Thiên lệch có thể xảy ra vì nhiều lý do. Việc phóng đại bởi một thước đo hoặc chỉ số thống kê có thể dẫn đến thiên lệch phóng đại (Rosenberg \& Houglet, 1974; Loughran \& Ritter, 1995). Khi các quan sát bị loại trừ khỏi mẫu do một quy tắc chọn lọc khác với việc lấy mẫu ngẫu nhiên, nó có thể tạo ra thiên lệch chọn mẫu. Thiên lệch sống sót là một ví dụ về thiên lệch chọn mẫu được thúc đẩy bởi sự loại trừ không tương xứng đối với các cổ phiếu đã bị hủy niêm yết theo thời gian (Waszczuk, 2014). Thiên lệch ươm tạo là một ví dụ về thiên lệch chọn mẫu do các quỹ mới có hiệu suất kém bị loại khỏi cơ sở dữ liệu hoặc không được đưa vào kịp thời (Wisen, 2002). Thiên lệch điền lùi phát sinh khi hiệu suất của quỹ không được công khai trong một khoảng thời gian ươm tạo nhưng sau đó được thêm vào cơ sở dữ liệu, có lẽ là sau một hiệu suất tốt (Jorion \& Schwarz, 2017). Thiên lệch nhìn trước xảy ra khi dữ liệu được sử dụng trong nghiên cứu được giả định là có sẵn cho công chúng tại một thời điểm cụ thể trong khi trên thực tế, dữ liệu đó chỉ có sẵn vào một thời điểm muộn hơn (Andrikopoulos et al., 2007). Việc bỏ qua lợi suất hủy niêm yết đối với một lượng lớn các công ty có thể dẫn đến thiên lệch hủy niêm yết (Shumway \& Warther, 1999; Waszczuk, 2014). Các thiên lệch có nhiều khả năng xảy ra đối với các công ty mới, công ty nhỏ, công ty nước ngoài, các công ty hủy niêm yết và các công ty có hiệu suất kém (Andrikopoulos et al., 2007; Landis \& Skouras, 2018). Thiên lệch có thể làm sai lệch kết quả nghiên cứu thực nghiệm và tạo ra các kết luận không đáng tin cậy. Một số nhà nghiên cứu cố gắng khắc phục thiên lệch sống sót bằng cách thêm các công ty đã hủy niêm yết trở lại vào mẫu của họ. Tuy nhiên, việc giảm thiểu các thiên lệch phụ thuộc rất nhiều vào những nỗ lực của các nhà cung cấp cơ sở dữ liệu trong việc cải thiện chất lượng dữ liệu và tính toàn vẹn của dữ liệu (data integrity).

\section{Sự Không nhất quán (Inconsistencies)}

Sự không nhất quán mô tả một tình huống mà dữ liệu không được xử lý một cách nhất quán với cùng một quy tắc, đặc biệt là trong cùng một cơ sở dữ liệu. Sự không nhất quán cần được điều tra thêm về các lỗi dữ liệu tiềm ẩn. Các vấn đề không nhất quán đã được báo cáo trong các cơ sở dữ liệu Compustat, Orbis và SDC (McElreath \& Wiggins, 1984, Faccio \& Masulis, 2005; Shi \& Zhang, 2011; Kalemli-Ozcan et al., 2019). Một ví dụ về sự không nhất quán là việc các khoản mục kế toán giống nhau được định nghĩa khác nhau trong các báo cáo tài chính khác nhau. Người ta đã báo cáo rằng trong Compustat ``chi phí khấu hao và phân bổ'' (depreciation and amortization expenses) được báo cáo trên báo cáo kết quả hoạt động kinh doanh và báo cáo lưu chuyển tiền tệ bao gồm các khoản mục khác nhau. Các sự kiện đặc biệt chẳng hạn như ``bán tài sản'' và ``ghi giảm hàng tồn kho'' (inventory write-offs) đã được điều chỉnh trong báo cáo lưu chuyển tiền tệ nhưng không phải trong bảng cân đối kế toán. Số tiền trình bày lại (restated amount) được phản ánh trong báo cáo lưu chuyển tiền tệ đã không được điều chỉnh trong bảng cân đối kế toán. Sự thay đổi về phân loại được phản ánh trong

% [PDF p.33 | in p.32]
bảng cân đối kế toán đã không được cập nhật trong báo cáo lưu chuyển tiền tệ (Shi \& Zhang, 2011). Hơn nữa, sự không nhất quán có thể xảy ra đối với các khoảng thời gian khác nhau, ví dụ, định nghĩa của Compustat về các khoản phải trả trong năm 2002 và 2003 là khác nhau (Shi \& Zhang, 2011). Các vấn đề không nhất quán cũng có thể xảy ra với các mã định danh độc quyền (proprietary identifiers). Một mã định danh độc quyền thường được nhận dạng là ID duy nhất để theo dõi một công ty; tuy nhiên, nó có thể không được sử dụng một cách nhất quán theo cách này. Các nhà nghiên cứu đã tiết lộ rằng mã BvD ID của một công ty có thể thay đổi do thay đổi địa chỉ, hình thức pháp lý (legal form), hoặc các hoạt động M\&A, điều này có thể gây ra vấn đề trong việc theo dõi các thay đổi của một công ty (Kalemli-Ozcan et al., 2019).

Các nhà nghiên cứu có thể giả định rằng các cơ sở dữ liệu từ cùng một nhà cung cấp có cùng một chính sách báo cáo. Nhưng Mathers and Giacomini (2016) đã nhắc nhở các nhà nghiên cứu rằng Compustat và Capital IQ đều từ S\&P Global Market Intelligence có chính sách mã hóa khác nhau về năm tài chính (Mathers \& Giacomini, 2016). I/B/E/S và Worldscope đều từ Thomson Reuters, trong khi 22\% sự chênh lệch giữa ngày công bố đã được báo cáo trong hai cơ sở dữ liệu này (Acker \& Duck, 2009). Thông tin từ các phần khác nhau của một cơ sở dữ liệu cũng có thể không nhất quán. Các nhà nghiên cứu đã báo cáo rằng trong SDC, dữ liệu ``phương thức thanh toán'' thu được từ phần mô tả và trường biến ``phương thức thanh toán'' có sự khác biệt thường xuyên (Faccio \& Masulis, 2005). Dữ liệu không nhất quán có thể tạo ra các vấn đề về khả năng so sánh (McElreath \& Wiggins, 1984). Nó có thể gây ra các vấn đề tương tự như lỗi dữ liệu và đôi khi đánh lừa các nhà nghiên cứu. Do các nhà nghiên cứu phải mất nhiều công sức để nhận thấy và giải quyết các điểm không nhất quán trong các cơ sở dữ liệu, vấn đề này làm suy yếu giá trị của các cơ sở dữ liệu thương mại.

\section{Dữ liệu Phần đầu Tĩnh (Static Header Data)}

Vấn đề dữ liệu phần đầu tĩnh (hoặc vấn đề về niên đại - vintage issue) xảy ra khi cơ sở dữ liệu chỉ cung cấp thông tin mới nhất có sẵn và thiếu các bản ghi chuỗi thời gian (Kalemli-Ozcan et al., 2019). Vấn đề dữ liệu phần đầu tĩnh xảy ra chủ yếu trong trường dữ liệu bao gồm ``tên công ty'', ``mã chứng khoán'' (ticker), ``địa chỉ'', ``trụ sở chính'', ``quyền sở hữu'' (ownership), ``sàn giao dịch chứng khoán'', và ``mã ngành'' (Ince \& Porter, 2006; Hines, 2016; Landis \& Skouras, 2018; Nobes \& Stadler, 2018; Kalemli-Ozcan et al., 2019). Các vấn đề như vậy đã được quan sát thấy trong Compustat, Orbis, Worldscope, và Datastream (Ince \& Porter, 2006; Hines, 2016; Landis \& Skouras, 2018; Nobes \& Stadler, 2018; Kalemli-Ozcan et al., 2019). Dữ liệu phần đầu tĩnh xóa đi hồ sơ theo dõi (track record) về các thay đổi trước đó, do đó nó hạn chế các nhà nghiên cứu đưa các biến dữ liệu này vào phân tích chuỗi thời gian. Vì dữ liệu có thể không đại diện chính xác cho thời điểm đó, khi sử dụng các trường dữ liệu phần đầu tĩnh để tạo mẫu, nó có thể dẫn đến mẫu không chính xác, thiếu sót, hoặc thiên lệch chọn mẫu (Landis \& Skouras, 2018).

\section{Sự Chuẩn hóa (Standardization)}

Sự chuẩn hóa và điều chỉnh làm cho dữ liệu có thể so sánh tốt hơn giữa các công ty, theo thời gian, và đặc biệt làm cho nó khả thi để so sánh các công ty từ các quốc gia khác nhau. Lợi ích rõ ràng của sự chuẩn hóa khiến nó trở thành điểm nhấn bán hàng của một số cơ sở dữ liệu như Compustat, Worldscope, và Orbis. Tuy nhiên, sự chuẩn hóa cũng gây ra các vấn đề về dữ liệu. Sự chuẩn hóa và điều chỉnh dữ liệu về ``khấu hao, cạn kiệt, và phân bổ'' (depreciation, depletion, and amortization) trong Compustat đã được phát hiện là ``báo cáo thấp đi 7.5\% so với giá vốn hàng bán trong hồ sơ 10-K'' và ``phóng đại tỷ suất lợi nhuận gộp trong hồ sơ 10-K thêm 14.3\%'' (Bostwick et al., 2016). Các nhà nghiên cứu cũng phát hiện ra sự chuẩn hóa trong Compustat ``không mang lại sự cải thiện nào cho các mô hình dự đoán phá sản và tác động tiêu cực đáng kể đến độ chính xác dự báo của mô hình Altman'' (Chychyla \& Kogan, 2014, tr. 1). Việc chuẩn hóa đối với

% [PDF p.34 | in p.33]
dữ liệu I/B/E/S được điều chỉnh theo chia tách cổ phiếu và làm tròn đến đồng xu gần nhất làm mất đi thông tin (Payne \& Thomas, 2003). Vì vậy, các nhà nghiên cứu nên cân nhắc kỹ hơn khi lựa chọn giữa tập dữ liệu đã chuẩn hóa và chưa chuẩn hóa. Cần phải nỗ lực hơn để hiểu quá trình chuẩn hóa và những thay đổi được tạo ra đối với các tập dữ liệu gốc do một quá trình như vậy.

\section{Những Thay đổi trong Dữ liệu Lịch sử (Changes in Historical Data)}

Dữ liệu lịch sử có thể được sửa đổi do việc sửa lỗi hoặc cập nhật từ sự trình bày lại. Nhưng những thay đổi có hệ thống trong dữ liệu lịch sử có thể là một vấn đề. Các nhà nghiên cứu sử dụng các đề xuất cổ phiếu của nhà phân tích I/B/E/S phát hiện ra các tệp dữ liệu trong cùng một khoảng thời gian được tải xuống tại các thời điểm khác nhau có những khác biệt đáng kể. Những khác biệt này bao gồm các mức đề xuất, việc bổ sung và xóa các bản ghi, xóa tên nhà phân tích, và các thay đổi trong các thuộc tính của dự báo thu nhập có sẵn trong mỗi phiên bản (Ljungqvist et al., 2008; Call et al., 2020). Những thay đổi tưởng chừng như không ngẫu nhiên này đã làm dấy lên những lo ngại về tính toàn vẹn của dữ liệu (Brown-Humes, 2006). Những thay đổi trong dữ liệu lịch sử có thể rất phổ biến do tính chất hay thay đổi của các cơ sở dữ liệu thương mại, đặc biệt đối với các nhà tổng hợp dữ liệu lấy dữ liệu từ các nguồn dữ liệu sơ cấp khác. Những thay đổi về các nhà cung cấp dữ liệu sơ cấp, những người đóng góp (contributors), và các điều khoản hợp đồng đều có thể dẫn đến những thay đổi trong cơ sở dữ liệu. Trong những trường hợp mà cơ sở dữ liệu trao cho những người đóng góp dữ liệu quyền kiểm soát trực tiếp đối với cơ sở dữ liệu, những thay đổi có nhiều khả năng xảy ra hơn và khó quản lý hơn.

\section{Thiếu tính Minh bạch (Lack of Transparency)}

Thiếu tính minh bạch là một vấn đề cơ bản đối với nhiều vấn đề dữ liệu khác. Nói chung, các nhà cung cấp cơ sở dữ liệu không minh bạch về quy trình quản lý và thu thập dữ liệu của họ, chứ chưa nói đến việc cảnh báo các nhà nghiên cứu về các vấn đề và thiên lệch dữ liệu tiềm ẩn (Annaert, et al., 2016). Các nhà nghiên cứu bày tỏ sự quan ngại về tính minh bạch của dữ liệu SDC về cách thức thu thập dữ liệu hoặc cách xác định các biến (Netter et al., 2011). Mặc dù các vấn đề dữ liệu thường không được giải thích công khai, đôi khi chúng được tiết lộ thông qua các cuộc trò chuyện riêng giữa các nhà nghiên cứu và các nhà cung cấp cơ sở dữ liệu. Khi bị đặt câu hỏi về những thay đổi trong dữ liệu lịch sử của các nhà phân tích cá nhân, I/B/E/S phản hồi rằng tên của các nhà phân tích cá nhân vẫn còn trong cơ sở dữ liệu, nhưng chúng không hiển thị trên các tệp mà giới học thuật nhìn thấy do nguồn cấp dữ liệu (data feed) không đầy đủ (Brown-Humes, 2006). Khi được hỏi về những thay đổi trong các dự báo thu nhập lịch sử, I/B/E/S giải thích rằng những điều chỉnh hồi tố có thể xảy ra do chia tách cổ phiếu, cổ tức bằng cổ phiếu, những điều chỉnh tiền tệ mặc định hoặc sửa lỗi và tiết lộ thêm rằng một số khác biệt xảy ra là do công ty môi giới duy trì quyền kiểm soát đối với việc phân phối những dự báo này và những người đăng ký (subscribers) học thuật thường chỉ có quyền truy cập vào một tập hợp con của tất cả các dự báo thu nhập được đóng góp vào cơ sở dữ liệu (Call et al., 2020). Bất chấp những lời giải thích từ các nhà cung cấp, việc thiếu tính minh bạch khiến các nhà nghiên cứu lo ngại và ảnh hưởng lớn đến niềm tin của các nhà nghiên cứu đối với các cơ sở dữ liệu thương mại (Call et al., 2020).

\section{Các Vấn đề về Thời gian Báo cáo (Reporting Time Issues)}

Thời gian báo cáo dữ liệu có thể là một vấn đề khi các nhà cung cấp cơ sở dữ liệu không minh bạch về thời gian báo cáo dữ liệu và lịch cập nhật của họ. Nó có thể ảnh hưởng đến nghiên cứu khi thời điểm mà dữ liệu có sẵn cho công chúng khác với thời điểm mà một nghiên cứu giả định nó là (McElreath \& Wiggins, 1984). Việc ghi nhận không đúng thời gian báo cáo có thể gây ra thiên lệch nhìn trước (Andrikopoulos, et al., 2007). Sự trễ báo cáo (reporting lag) hoặc sự chậm trễ là một vấn đề khác đôi khi là không thể tránh khỏi. Trong trường hợp dữ liệu được tổng hợp từ các nhà cung cấp khác nhau, sự

% [PDF p.35 | in p.34]
trễ báo cáo có thể gây ra bởi các điều khoản hợp đồng cấm tiết lộ (embargo) cụ thể. Trong trường hợp dữ liệu có nguồn gốc từ các quốc gia khác nhau, sự trễ báo cáo cũng có thể do các thực tiễn tuân thủ báo cáo khác nhau ở các quốc gia khác nhau (Green, 2003). Độ trễ báo cáo trong Orbis được phát hiện là khoảng 2 năm tính trung bình và thay đổi theo từng quốc gia và từng sản phẩm dữ liệu (Kalemli-Ozcan et al. 2019). Độ trễ báo cáo do sự ươm tạo (khi các quỹ mới có hiệu suất kém không được đưa vào cơ sở dữ liệu nhanh chóng như các quỹ mới có hiệu suất vượt trội) có thể dẫn đến thiên lệch chọn mẫu (Wisen, 2002).

\section{Sự Lạm dụng Dữ liệu (Misuse of Data)}

Các vấn đề dữ liệu không chỉ giới hạn ở bản thân dữ liệu, đôi khi các nhà nghiên cứu có thể sử dụng dữ liệu một cách không đúng đắn như là các biến đại diện (proxies) hoặc các thước đo. Ví dụ, việc ghép (nhân gộp) lợi suất hàng ngày (compounding) của chỉ số tỷ trọng đều (equal-weighted index) CRSP để tính lợi suất hàng tháng có thể dẫn đến những thiên lệch lớn (Canina et al., 1998). Các hệ số điều chỉnh (adjustment factors) đối với giá cổ phiếu của CRSP đã bao gồm các hiệu ứng từ các sự kiện cổ tức bằng tài sản, chia tách công ty (spin-off), và chào bán quyền mua cổ phần (rights offering), do đó các nhà nghiên cứu chỉ được sử dụng các hệ số điều chỉnh của CRSP để xử lý (accommodate) các sự kiện thiếu tính thực chất kinh tế như vậy, nếu không, nó có thể tạo ra các quan sát mẫu sai sót và các kết quả sai lệch (Francis et al., 2016). Mills et al. (2003) cảnh báo các nhà nghiên cứu cần đặc biệt thận trọng khi sử dụng dữ liệu lỗ hoạt động ròng của Compustat làm chỉ báo cho các vị thế lỗ thuế ở Hoa Kỳ của một công ty, đặc biệt là khi bối cảnh nghiên cứu liên quan đến các công ty có hoạt động ở nước ngoài hoặc hoạt động mua lại của công ty. Hơn nữa, vì Compustat chỉ bao phủ các công ty đại chúng trong một ngành, nó là một đại diện kém cho mức độ tập trung ngành thực tế. Việc xây dựng các thước đo tập trung ngành bằng cách sử dụng dữ liệu Compustat có thể dẫn đến kết luận không chính xác (Ali et al., 2008; Keil, 2017). Những thông báo mua lại cổ phần trong SDC là những công cụ dự báo kém đối với số lượng cổ phiếu sẽ được công ty mua lại, bởi vì dữ liệu không đầy đủ (Banyi et al., 2008).

Thông qua nghiên cứu, chúng tôi thấy rằng một số vấn đề chất lượng dữ liệu có thể được giảm nhẹ thông qua sự cải thiện trong các phương pháp kiểm định thống kê, các kỹ thuật lập hồ sơ và khai phá dữ liệu (data profiling and mining), cùng với việc áp dụng các chính sách và hệ thống báo cáo dữ liệu mới. Tuy nhiên, sự cải thiện chất lượng dữ liệu tổng thể phần lớn phụ thuộc vào các hoạt động kiểm soát chất lượng minh bạch của các nhà cung cấp dữ liệu và sự tham gia cởi mở của họ với cộng đồng nghiên cứu.

\section{Hàm ý đối với Tham khảo Kinh doanh và Tư vấn Nghiên cứu}

Giúp người dùng thư viện và các nhà nghiên cứu hiểu chất lượng thông tin là một phần thiết yếu của các dịch vụ tham khảo thư viện. Các thủ thư đã đặc biệt chú ý đến các khía cạnh chất lượng thông tin như độ tin cậy, tính hợp lệ (validity), độ chính xác, tính thẩm quyền, tính kịp thời, và các thiên lệch (ACRL, 2000). Chúng tôi cũng đã phát triển các phương pháp thực tiễn như bộ tiêu chí CRAAP (Tính cập nhật - Currency, Tính phù hợp - Relevance, Thẩm quyền - Authority, Độ chính xác - Accuracy, và Mục đích - Purpose) để tạo điều kiện cho giáo dục về hiểu biết thông tin (Blakeslee, 2004). Tuy nhiên, nghiên cứu này nhận thấy các vấn đề chất lượng dữ liệu có thể phức tạp và đáng lo ngại hơn nhiều. Những thủ thư kinh doanh cần phải nhận thức được các vấn đề chất lượng dữ liệu này và nâng cao nhận thức của các nhà nghiên cứu về các vấn đề này thông qua tham khảo kinh doanh và tư vấn. Dưới đây là một số lời khuyên thực tế mà chúng ta có thể cung cấp cho các nhà nghiên cứu.

% [PDF p.36 | in p.35]
\begin{itemize}
  \item Các nhà nghiên cứu không nên coi chất lượng của các cơ sở dữ liệu thương mại danh tiếng là điều hiển nhiên. Họ luôn nên đánh giá chất lượng dữ liệu và kiểm tra độ chính xác và tính đầy đủ của nó. Họ cần cẩn trọng với các dự án nghiên cứu những khái niệm kế toán hoặc các giao dịch kinh doanh phức tạp, hoặc các dự án liên quan đến các công ty mới, công ty nhỏ, công ty nước ngoài, hoặc công ty hủy niêm yết.
  \item Các nhà nghiên cứu không nên chỉ dựa vào một nguồn dữ liệu duy nhất. Nếu có thể, họ nên tham khảo nhiều nguồn, đặc biệt là các nguồn dữ liệu gốc. Họ nên chú ý đến các định nghĩa và chính sách mã hóa khác nhau của các cơ sở dữ liệu khác nhau và coi sự chênh lệch giữa các cơ sở dữ liệu như một cơ hội để xác định các giá trị bị thiếu hoặc lỗi dữ liệu. Họ nên cẩn trọng vì ``hiệu ứng cơ sở dữ liệu'' và kiểm chứng các lý thuyết của họ thông qua nhiều nguồn dữ liệu.
  \item Các nhà nghiên cứu nên đọc cẩn thận các cẩm nang dữ liệu và hiểu cách dữ liệu được định nghĩa hoặc tính toán, đặc biệt đối với dữ liệu đã được chuẩn hóa hoặc điều chỉnh. Họ nên nhận thức được sự không nhất quán của cơ sở dữ liệu trong việc xử lý các định nghĩa và những điều chỉnh như vậy trên các báo cáo tài chính, các khoảng thời gian lịch sử, và xuyên suốt cơ sở dữ liệu.
  \item Các nhà nghiên cứu nên thận trọng trong việc sử dụng cơ sở dữ liệu làm công cụ sàng lọc để xác định các mẫu dữ liệu, đặc biệt là sử dụng trường dữ liệu phần đầu tĩnh như một biến để sàng lọc các mẫu dữ liệu. Họ cần đánh giá xem liệu dữ liệu mẫu có phải là các biến đại diện hoặc thước đo thích hợp hay không. Dữ liệu không đầy đủ và các phân loại không chính xác trong cơ sở dữ liệu có thể gây ra các bản ghi bị thiếu hoặc không chính xác trong các mẫu dữ liệu hoặc thậm chí tạo ra thiên lệch chọn mẫu và kết quả sai lệch.
  \item Các nhà nghiên cứu không nên bỏ qua các thiên lệch trong cơ sở dữ liệu. Thay vào đó, họ cần tìm kiếm các thủ tục thích hợp để giảm thiểu các thiên lệch và giải thích rõ ràng về tác động của các thiên lệch tiềm ẩn cũng như những hạn chế của nghiên cứu của họ.
  \item Các nhà nghiên cứu không nên coi việc thu thập dữ liệu là giao dịch một lần (one-time transaction). Thay vào đó, họ cần hiểu tính chất hay thay đổi của một cơ sở dữ liệu nghiên cứu, bảo tồn dữ liệu ở các thời điểm khác nhau, và để lại một bằng chứng bằng văn bản (paper trail) về truy cập và sử dụng dữ liệu.
  \item Các nhà nghiên cứu cần nhận thức được thời gian báo cáo, độ trễ báo cáo, lịch cập nhật, hoặc khoảng thời gian cấm tiết lộ (embargo) của các nguồn dữ liệu. Họ có thể cần điều chỉnh cách thu thập dữ liệu hoặc các thông lệ cập nhật của họ cho phù hợp.
  \item Các nhà nghiên cứu nên giữ liên lạc cởi mở với các nhà cung cấp về các vấn đề dữ liệu. Việc trao đổi của họ sẽ cho phép các nhà cung cấp xác định những vấn đề tương tự hoặc khởi xướng các dự án để tạo ra những thay đổi ở quy mô lớn.
\end{itemize}
\section{Kết luận}

Bài viết này cung cấp một bài tổng quan tài liệu về các tài liệu kinh doanh và tài chính đề cập đến các vấn đề về chất lượng dữ liệu và bao gồm các cơ sở dữ liệu như CRSP, Compustat, Capital IQ, I/B/E/S, Datastream, Worldscope, SDC, và BvD Orbis. Sự phân tích tổng hợp (synthetic analysis) trên các tài liệu kinh doanh đã xác định được 11 loại vấn đề chất lượng dữ liệu phổ biến, bao gồm giá trị bị thiếu, lỗi dữ liệu, sự chênh lệch, các thiên lệch, sự không nhất quán, dữ liệu phần đầu tĩnh, sự chuẩn hóa, những thay đổi trong dữ liệu lịch sử, thiếu tính minh bạch, các vấn đề về thời gian báo cáo và sự lạm dụng dữ liệu. Các vấn đề dữ liệu này có thể đưa vào các lỗi và thiên lệch vào trong nghiên cứu thực nghiệm bằng nhiều cách, làm nhiễu phân tích thống kê, và bóp méo kết quả nghiên cứu. Nhiều nhà nghiên cứu đã bác bỏ

% [PDF p.37 | in p.36]
các kết quả nghiên cứu trước đó sau khi sửa chữa các vấn đề dữ liệu cụ thể. Bất chấp sự phổ biến của các vấn đề về chất lượng dữ liệu, các cơ sở dữ liệu này được các nhà nghiên cứu học thuật sử dụng rộng rãi để thực hiện các nghiên cứu thực nghiệm trong các lĩnh vực kế toán, tài chính, kinh tế, toán học và thống kê, bởi thị trường thương mại cho mục đích kiểm thử ngược (backtesting) và tính toán mô hình, và bởi các cơ quan chính phủ phục vụ phân tích tài chính và kinh tế. Nghiên cứu học thuật từ lâu đã được tin cậy như một cách thức đáng tin cậy để tạo ra tri thức và đạt được những tiến bộ về mặt khoa học và lý thuyết trong các lĩnh vực liên quan. Việc ra quyết định dựa trên bằng chứng và dựa trên dữ liệu đã được hoan nghênh rộng rãi như một thực tiễn ra quyết định đáng tin cậy hơn. Các vấn đề chất lượng dữ liệu không chỉ làm suy yếu giá trị của nghiên cứu học thuật, đánh lừa các quyết định kinh doanh, gây nhiễu (confound) các chính sách của chính phủ mà còn làm tổn hại đến niềm tin của công chúng vào tri thức. Các thủ thư đã đóng một vai trò quan trọng trong việc hỗ trợ cộng đồng nghiên cứu truy cập các nguồn dữ liệu và giáo dục cho các nhà nghiên cứu. Hy vọng rằng, bài viết này sẽ giúp tạo điều kiện thuận lợi cho sự giao tiếp của thủ thư với các nhà nghiên cứu về các vấn đề chất lượng dữ liệu và sẽ nâng cao nhận thức của cộng đồng nghiên cứu về vấn đề này.

\section{Phụ lục I Danh sách các Bài báo được Đánh giá (List of Reviewed Articles)}

{\small
\begin{longtable}{>{\raggedright\arraybackslash}p{0.241\linewidth}>{\raggedright\arraybackslash}p{0.236\linewidth}>{\raggedright\arraybackslash}p{0.462\linewidth}}
\toprule
\textbf{Bài báo (Tác giả-Năm)} & \textbf{Tên Tạp chí} & \textbf{Tóm tắt (cung cấp cho cơ sở dữ liệu có trên 5 bài báo được trích dẫn)} \\
\midrule
\endhead
\textbf{Center for Research in Security Prices (CRSP)} &  &  \\
\addlinespace[2pt]
Rosenberg and Houglet (1974) & The Journal of Finance & Phạm vi Thời gian: Số lượng Bài báo\newline{}1970-1999: 12\newline{}2000-2009: 4\newline{}2010-2019: 5\newline{}Phạm vi Tạp chí: Số lượng Bài báo\newline{}The Journal of Finance: 9\newline{}Social Science Research Network (SSRN): 4\newline{}Journal of Financial and Quantitative Analysis: 2\newline{}Journal of Financial Economics: 1\newline{}The Accounting Review: 1\newline{}The Review of Financial Studies: 1\newline{}Journal of Accounting and Economics: 1\newline{}Working Paper: 1\newline{}Manuscript: 1\newline{}Tổng cộng: 21 \\
\addlinespace[2pt]
Beedles and Simkowitz (1978) & The Journal of Finance &  \\
\addlinespace[2pt]
Bennin (1980) & The Journal of Finance &  \\
\addlinespace[2pt]
Grinblatt et al. (1984) & Journal of Financial Economics &  \\
\addlinespace[2pt]
Sarig and Warga (1989) & Journal of Financial and Quantitative Analysis &  \\
\addlinespace[2pt]
Guenther and Rosman (1994) & Journal of Accounting and Economics &  \\
\addlinespace[2pt]
Courtenay and Keller (1994) & The Accounting Review &  \\
\addlinespace[2pt]
Loughran and Ritter (1995) & The Journal of Finance &  \\
\addlinespace[2pt]
Kahle and Walkling (1996) & Journal of Financial and Quantitative Analysis &  \\
\addlinespace[2pt]
Shumway (1997) & The Journal of Finance &  \\
\addlinespace[2pt]
Canina (1998) & The Journal of Finance &  \\
\addlinespace[2pt]
Shumway and Warther (1999) & The Journal of Finance &  \\
\addlinespace[2pt]
Elton et al. (2001) & The Journal of Finance &  \\
\addlinespace[2pt]
Wisen (2002) & SSRN &  \\
\addlinespace[2pt]
Evans (2007) & Manuscript &  \\
\addlinespace[2pt]
Yan (2007) & SSRN &  \\
\addlinespace[2pt]
Evans (2010) & The Journal of Finance &  \\
\addlinespace[2pt]
Schwarz and Potter (2016) & The Review of Financial Studies &  \\
\addlinespace[2pt]
Francis et al. (2016) & SSRN &  \\
\addlinespace[2pt]
Jorion and Schwarz (2017) & SSRN &  \\
\addlinespace[2pt]
Tobek and Hronec (2018) & IES Working Paper &  \\
\addlinespace[2pt]
\textbf{Compustat} &  &  \\
\addlinespace[2pt]
San Miguel (1977) & The Accounting Review & Phạm vi Thời gian: Số lượng Bài báo\newline{}1970-1999: 6\newline{}2000-2009: 6 \\
\addlinespace[2pt]
Ball (1979) & The Journal of Finance &  \\
\addlinespace[2pt]
McElreath and Wiggins (1984) & Financial Analysts Journal &  \\
\addlinespace[2pt]
\bottomrule
\end{longtable}
}

\emph{(Bảng Phụ lục I tiếp tục ở trang in 37, ngoài phạm vi đọc của WQU; ô tóm tắt của Compustat cũng tiếp tục ở đó nên số liệu trong ô này chưa phải tổng cuối cùng.)}
